\chapter{Benchmark、profiling 与优化方法}

\begin{introduction}
  \item 把性能问题写成可反驳假设
  \item 区分微基准、端到端基准与 profiler
  \item 正确处理预热、噪声、优化器和统计摘要
  \item 用性能预算而不是“越快越好”指导取舍
\end{introduction}

\section{优化循环}

可重复的优化循环是：定义用户可见指标，记录基线，定位热点，提出单一假设，做最小改动，重新测量，并验证正确性没有回归。若结果落在噪声范围内，结论是“证据不足”，不是挑最好的一次宣布成功。

吞吐量适合批量历史数据；平均延迟会掩盖尾部；交易系统常关注 \(p_{50}\)、\(p_{99}\)、\(p_{99.9}\) 和最大值，同时记录负载和丢弃策略。任何百分位都必须说明样本量、采样边界和计时器。

\section{微基准的陷阱}

优化器可能删除未使用结果，因此基准应消费 checksum；首次运行包含页面映射和缓存冷启动；CPU 动态频率、后台任务和安全软件带来噪声；计时器调用本身也有成本。微基准适合隔离机制，不代表整套系统收益。

本书的布局程序刻意同时输出结果一致性和耗时。严谨版本应增加多次迭代、随机化顺序、预热、统计摘要，并在固定硬件上运行。Google Benchmark 等框架能处理部分细节，但不能替你定义有意义工作负载。

\section{profiler 回答“时间在哪里”}

采样 profiler 以较低扰动发现 CPU 热点和调用栈；instrumentation 可提供精确事件但扰动更大；硬件计数器观察周期、指令、分支、缓存未命中。先用低成本方法定位，再针对假设选择证据。

如果函数占总时间 \(5\%\)，即使把它变为零，整体最多约提升 \(1/(1-0.05)\approx1.053\) 倍。Amdahl 定律提醒我们不要优化醒目的小函数而忽略真正瓶颈。

\section{端到端预算}

回测性能可以分解为读取、解析、事件分发、策略、撮合、簿记和统计。实时路径还包括网络、排队、序列化与系统调用。为每段定义预算和观测点，才能判断局部微基准是否移动整体指标。

\section{性能与可维护性}

若两种实现差异小于噪声，选择更清楚者。若优化显著但复杂，应记录适用工作负载、基准命令和回退方案，并用回归基准保护。不要把硬件相关技巧扩散到业务层；将其封装在有普通参考实现的深模块里。

\begin{interviewcheck}
为什么 checksum 能阻止部分死代码删除？微基准与端到端基准回答什么不同问题？如何解释 \(p_{99}\)？Amdahl 定律如何改变优化优先级？
\end{interviewcheck}

\section{练习}

\begin{enumerate}
  \item 为布局基准增加 20 次测量，报告中位数、最小值和四分位距。
  \item 用 profiler 检查回测样例，提出一个可被数据推翻的瓶颈假设。
  \item 定义每秒事件数与 \(p_{99}\) 延迟的双重性能预算。
  \item 分析一个优化如果使代码复杂一倍但只提升 \(2\%\)，需要哪些业务证据才值得保留。
\end{enumerate}

\section{本章小结}

性能工程是一种实验方法。指标、基线、热点、假设、改动和复测必须形成闭环；微基准提供局部证据，profiler 定位系统成本，端到端预算决定优化是否有价值。
